% Thesis-ready section for Chapter 4.
% Source guardrails: results/six_method_v1/SIX_METHOD_REVIEW.md and SIX_METHOD_AUDIT.json.

\section{Giới hạn của phân tích thống kê và phạm vi diễn giải}
\label{sec:statistical-limitations}

Các kết quả của luận văn được xây dựng trên một giao thức khóa gồm năm kích thước STAR-RIS, tám seed huấn luyện cho mỗi thuật toán học tăng cường sâu và 1.000 kịch bản kiểm thử cho mỗi kích thước. Thiết kế này cho phép so sánh có kiểm soát giữa các phương pháp, nhưng không loại bỏ hoàn toàn bất định do quá trình huấn luyện, mô hình kênh và nền tảng phần cứng. Vì vậy, các kết luận dưới đây chỉ áp dụng cho phạm vi thí nghiệm đã công bố.

\subsection{Đơn vị bất định không hoàn toàn giống nhau giữa hai nhóm phương pháp}

Đối với TD3, DDPG và PPO, các thống kê tổng hợp trước hết được tính trên từng seed rồi mới mô tả sự biến thiên giữa tám seed. Ngược lại, AO-SCA, AO-Grid và AnalyticalRIS là các phương pháp xác định hoặc gần xác định trong giao thức hiện tại, nên độ lệch chuẩn của chúng chủ yếu phản ánh sự khác nhau giữa 1.000 kịch bản kênh. Do đó, khoảng tin cậy của nhóm DRL và nhóm tối ưu truyền thống không đại diện cho cùng một nguồn bất định. Luận văn sử dụng các khoảng này để mô tả kết quả trong từng nhóm, nhưng không diễn giải độ rộng của chúng như một phép so sánh trực tiếp về độ ổn định giữa hai loại phương pháp.

\subsection{Kiểm định ghép cặp chưa phải là mô hình phân cấp đầy đủ}

Các kiểm định ghép cặp sử dụng cùng 1.000 kịch bản kiểm thử. Với các phương pháp DRL, kết quả của tám seed được trung bình theo từng kịch bản trước khi thực hiện kiểm định. Cách làm này tránh xem tám seed như tám bản sao độc lập của cùng một kịch bản, nhưng nó kiểm định sự khác biệt theo kịch bản của một tập hữu hạn các chính sách đã huấn luyện. Phân tích chưa xây dựng mô hình phân cấp đồng thời cho hai nguồn biến thiên là seed huấn luyện và kịch bản kiểm thử. Vì vậy, các giá trị $p$ không được xem là bằng chứng tuyệt đối rằng kết luận sẽ giữ nguyên với mọi seed hoặc mọi phân bố triển khai khác.

\subsection{Hiệu chỉnh nhiều so sánh và hiện tượng tràn số của giá trị $p$}

Bảng kiểm định sử dụng hiệu chỉnh Holm cho họ các so sánh được báo cáo. Một số giá trị $p$ rất nhỏ được phần mềm lưu thành 0 do giới hạn số học dấu phẩy động. Trong luận văn, các trường hợp này phải được viết dưới dạng ``nhỏ hơn ngưỡng biểu diễn'', chẳng hạn $p<10^{-300}$, thay vì khẳng định $p=0$. Bên cạnh giá trị $p$, luận văn đồng thời báo cáo chênh lệch trung bình và Cohen's $d_z$ để thể hiện độ lớn thực tế của hiệu ứng.

\subsection{Hiệu ứng rất lớn có thể xuất phát từ phương sai rất nhỏ}

Một số baseline như AO-Grid và AnalyticalRIS tạo ra kết quả gần như không thay đổi giữa các kịch bản trong thiết lập hiện tại. Khi độ lệch chuẩn của chênh lệch rất nhỏ, Cohen's $d_z$ có thể đạt giá trị rất lớn. Những giá trị này không nên được hiểu theo thang diễn giải thông thường một cách cơ học. Chênh lệch tuyệt đối về tổng tốc độ, QoS và độ trễ vẫn là các đại lượng chính để đánh giá ý nghĩa thực tiễn.

\subsection{Các chỉ số QoS là đại lượng bị chặn}

Tỷ lệ người dùng đạt QoS và xác suất toàn bộ người dùng đạt QoS nằm trong đoạn $[0,1]$. Khoảng tin cậy Student-$t$ không bị chặn có thể vượt nhẹ khỏi miền vật lý khi trung bình gần 0 hoặc 1. Dữ liệu thô được giữ nguyên để bảo đảm khả năng tái lập; khi trình bày trực quan, luận văn giới hạn trục trong $[0,1]$ và không sử dụng phần vượt miền vật lý để đưa ra kết luận.

\subsection{Giới hạn của phép đo độ trễ}

Độ trễ được đo bằng 100 lần lặp cho mỗi phương pháp và mỗi $N$ trên cùng một CPU runner đơn luồng. Thiết kế này giúp so sánh tương đối trong cùng môi trường, nhưng không đại diện cho mọi CPU, GPU, hệ điều hành hoặc thư viện số. Các tỷ lệ tăng tốc chỉ được phát biểu cho nền tảng đã công bố. AnalyticalRIS có độ trễ nhỏ nhất nhưng không đáp ứng QoS; do đó, độ trễ không thể được diễn giải tách rời chất lượng nghiệm.

\subsection{Giới hạn của tính công bằng siêu tham số}

Ba thuật toán TD3, DDPG và PPO dùng cùng ScenarioBank, số tương tác, hàm thưởng QoS, quy tắc chọn checkpoint và tập kiểm thử. Tuy nhiên, thí nghiệm không thực hiện một quá trình tìm kiếm siêu tham số có ngân sách bằng nhau cho mọi thuật toán. TD3 sử dụng Layer Normalization, Huber loss, cắt gradient và chuẩn hóa nhiễu theo chiều action, trong khi DDPG và PPO dùng các thiết lập đơn giản hơn. Vì vậy, kết luận phù hợp là TD3 đạt kết quả tốt nhất trong ba implementation được đánh giá theo giao thức đã công bố; luận văn không khẳng định TD3 luôn vượt mọi cấu hình DDPG hoặc PPO có thể xây dựng.

\subsection{Giới hạn của mô hình và khả năng khái quát}

Các kịch bản kiểm thử được sinh từ mô hình kênh tổng hợp với CSI hoàn hảo, người dùng SISO, vị trí miền truyền và phản xạ cố định, SIC lý tưởng và STAR-RIS thụ động với pha truyền và phản xạ độc lập. Kết quả chưa chứng minh hiệu năng dưới sai số CSI, tương quan không gian thực tế, phần cứng lượng tử hóa, ghép pha vật lý, di động người dùng hoặc dữ liệu đo ngoài hiện trường. AO-SCA cũng chỉ là baseline lặp cục bộ phụ thuộc khởi tạo, không phải nghiệm tối ưu toàn cục hay cận trên lý thuyết.

Tóm lại, kết quả cho phép kết luận rằng TD3 là phương pháp DRL ổn định và có khả năng mở rộng tốt nhất trong ba thuật toán học được đánh giá, đồng thời tạo ra sự cân bằng mạnh giữa tổng tốc độ, độ tin cậy QoS và độ trễ. Kết luận này phải luôn đi kèm phạm vi của giao thức khóa, tám seed huấn luyện, mô hình kênh tổng hợp và nền tảng CPU đã sử dụng.